% Mathematical Language style and canonical notation
% (MATH000) Mathematical Language - Notation and Conventions
%
% Design baseline: Mathematics for Science Quiz, with a lighter reference-document
% shell. Institution colors are reserved for institution names only.

\usepackage[a4paper,margin=1in,headheight=14pt,headsep=18pt]{geometry}
\usepackage[T1]{fontenc}
\usepackage{lmodern}
\usepackage{microtype}
\usepackage[intlimits]{amsmath}
\usepackage{amssymb,amsthm}
\usepackage{mathtools}
\usepackage{bbm}
\usepackage{mathrsfs}
\usepackage{xcolor}
\usepackage{array}
\usepackage{booktabs}
\usepackage{longtable}
\usepackage{enumitem}
\usepackage{titlesec}
\usepackage{tocloft}
\usepackage{fancyhdr}
\usepackage{needspace}
\usepackage{etoolbox}
\usepackage[hidelinks]{hyperref}

% Institution colors
\definecolor{POSTECHRed}{RGB}{166,25,85}
\newcommand{\institutiontext}[2]{\texorpdfstring{\textcolor{#1}{#2}}{#2}}
\DeclareRobustCommand{\POSTECH}{\institutiontext{POSTECHRed}{POSTECH}}
\DeclareRobustCommand{\PohangUniversityOfScienceAndTechnology}{%
  \institutiontext{POSTECHRed}{Pohang University of Science and Technology}%
}

% Number systems and common objects
\newcommand{\N}{\mathbb{N}}
\newcommand{\Nzero}{\mathbb{N}_0}
\newcommand{\Z}{\mathbb{Z}}
\newcommand{\Q}{\mathbb{Q}}
\newcommand{\R}{\mathbb{R}}
\newcommand{\C}{\mathbb{C}}
\newcommand{\F}{\mathbb{F}}
\newcommand{\Hquat}{\mathbb{H}}
\newcommand{\Prob}{\mathbb{P}}
\newcommand{\Exp}{\mathbb{E}}
\newcommand{\indicator}[1]{\mathbbm{1}_{#1}}
\newcommand{\Pow}{\mathcal{P}}

% Upright differential
\newcommand{\dd}{\mathop{}\!\mathrm{d}}

% Logarithms and inverse trigonometric functions
\DeclareMathOperator{\lb}{lb}
\DeclareMathOperator{\arcsec}{arcsec}
\DeclareMathOperator{\arccsc}{arccsc}
\DeclareMathOperator{\arccot}{arccot}

% Named operators
\DeclareMathOperator{\sgn}{sgn}
\DeclareMathOperator{\dom}{dom}
\DeclareMathOperator{\codom}{codom}
\DeclareMathOperator{\id}{id}
\DeclareMathOperator{\rk}{rk}
\DeclareMathOperator{\tr}{tr}
\DeclareMathOperator{\adj}{adj}
\DeclareMathOperator{\Col}{Col}
\DeclareMathOperator{\Row}{Row}
\DeclareMathOperator{\Null}{Null}
\DeclareMathOperator{\Span}{Span}
\DeclareMathOperator{\im}{im}
\DeclareMathOperator{\RePart}{Re}
\DeclareMathOperator{\ImPart}{Im}
\DeclareMathOperator{\diag}{diag}
\DeclareMathOperator{\Jac}{Jac}
\DeclareMathOperator{\Hess}{Hess}
\DeclareMathOperator{\Res}{Res}
\DeclareMathOperator{\spec}{spec}
\DeclareMathOperator{\ord}{ord}
\DeclareMathOperator{\rad}{rad}
\DeclareMathOperator{\Aut}{Aut}
\DeclareMathOperator{\GL}{GL}
\DeclareMathOperator{\SL}{SL}
\DeclareMathOperator{\Gal}{Gal}
\DeclareMathOperator{\charac}{char}
\DeclareMathOperator{\trdeg}{trdeg}
\DeclareMathOperator{\Var}{Var}
\DeclareMathOperator{\Cov}{Cov}
\DeclareMathOperator{\Bin}{Bin}
\DeclareMathOperator{\Poi}{Pois}
\newcommand{\Normal}{\mathcal{N}}
\DeclareMathOperator{\dist}{dist}
\DeclareMathOperator{\conv}{conv}
\DeclareMathOperator{\area}{area}
\DeclareMathOperator{\vol}{vol}
\DeclareMathOperator{\lengthop}{length}
\DeclareMathOperator{\interior}{int}
\DeclareMathOperator{\closure}{cl}
\DeclareMathOperator{\lcmop}{lcm}

% Common vector/matrix symbols
\newcommand{\ones}{\mathbf{1}}

% Convergence conventions
\newcommand{\ptwto}{\to}
\newcommand{\unifto}{\rightrightarrows}
\newcommand{\ptwlabelto}{\xrightarrow{\mathrm{ptw}}}
\newcommand{\uniflabelto}{\xrightarrow{\mathrm{unif}}}
\newcommand{\aeto}{\xrightarrow{\mathrm{a.e.}}}
\newcommand{\asto}{\xrightarrow{\mathrm{a.s.}}}
\newcommand{\probto}{\xrightarrow{\mathbb{P}}}
\newcommand{\distto}{\xrightarrow{\mathrm{d}}}
\newcommand{\Lpto}[1]{\xrightarrow{L^{#1}}}
\newcommand{\indep}{\mathrel{\perp\!\!\!\perp}}

% Proof-ending conventions
\newcommand{\ordinaryqed}{\ensuremath{\square}}
\newcommand{\majorqed}{\ensuremath{\blacksquare}}
\newcommand{\exceptionalqed}{\textsc{Q.E.D.}}

% Page typography
\setlength{\parskip}{0pt}
\setlength{\parindent}{1.5em}
\linespread{1.0}

% Table columns and spacing
\newcolumntype{L}[1]{>{\raggedright\arraybackslash}p{#1}}
\setlength{\LTpre}{0.5em}
\setlength{\LTpost}{0.62em}

\newenvironment{notationtable}{%
  \par\vspace{0.05em}%
  \small
  \renewcommand{\arraystretch}{1.18}%
  \begin{longtable}{@{}L{0.27\textwidth}L{0.65\textwidth}@{}}%
  \toprule
  Notation & Meaning \\
  \midrule
  \endfirsthead
  \toprule
  Notation & Meaning \\
  \midrule
  \endhead
  \bottomrule
  \endlastfoot
}{%
  \end{longtable}%
}

\newenvironment{ruletable}{%
  \par\vspace{0.05em}%
  \small
  \renewcommand{\arraystretch}{1.18}%
  \begin{longtable}{@{}L{0.24\textwidth}L{0.68\textwidth}@{}}%
  \toprule
  Item & Convention \\
  \midrule
  \endfirsthead
  \toprule
  Item & Convention \\
  \midrule
  \endhead
  \bottomrule
  \endlastfoot
}{%
  \end{longtable}%
}

% Contents
\setlength{\cftbeforepartskip}{0.52em}
\setlength{\cftbeforesecskip}{0.11em}
\renewcommand{\cftpartfont}{\normalsize\bfseries}
\renewcommand{\cftpartpagefont}{\normalsize}
\renewcommand{\cftsecfont}{\small}
\renewcommand{\cftsecpagefont}{\small}
\setlength{\cftsecnumwidth}{2.0em}

% Lightweight Parts
\newcounter{mathpart}
\newcommand{\mathpart}[1]{%
  \Needspace{11\baselineskip}%
  \refstepcounter{mathpart}%
  \addcontentsline{toc}{part}{Part \Roman{mathpart}. #1}%
  \vspace{1.15em}%
  {\noindent\large\bfseries Part \Roman{mathpart}. #1\par}%
  \vspace{0.68em}%
}

% Section hierarchy
\titleformat{\section}
  {\large\bfseries}
  {\thesection}
  {0.72em}
  {}
\titlespacing*{\section}{0pt}{1.40em}{0.62em}
\pretocmd{\section}{\Needspace{8\baselineskip}}{}{}

% Subsections are reserved for future use.
\titleformat{\subsection}
  {\normalsize\bfseries}
  {\thesubsection}
  {0.68em}
  {}
\titlespacing*{\subsection}{0pt}{1.1em}{0.45em}

% Running identity
\pagestyle{fancy}
\fancyhf{}
\fancyhead[L]{\small\textit{(MATH000) Mathematical Language}}
\fancyhead[R]{\small\textit{Notation and Conventions}}
\fancyfoot[C]{\small\thepage}
\renewcommand{\headrulewidth}{0.35pt}
\renewcommand{\footrulewidth}{0pt}

\setlist[itemize]{leftmargin=1.4em,itemsep=0.15em,topsep=0.25em}
\setlist[enumerate]{leftmargin=1.6em,itemsep=0.15em,topsep=0.25em}

\hypersetup{
  pdftitle={Mathematical Language: Notation and Conventions},
  pdfauthor={Jungwoo Kim},
  pdfsubject={Mathematical notation and writing conventions},
  pdfcreator={LaTeX}
}
